\chapter{Servers}\label{ch:nw:servers}

The servers that trading firms buy for the hot path are chosen for one core's speed, not for the number of cores. A mainstream
server processor with 128 cores runs any one of them at up to 4.1~GHz; a server processor with 16 cores runs one at up to 5~GHz;
and specialist vendors sell machines whose processor has been tuned to run all 24 of its cores at 4.8~GHz, cooled by liquid,
powered by a 2\,000-watt supply and warranted by the vendor rather than the chip maker. A trading path that spends most of its time
computing gets faster almost in proportion to the clock; one that spends it waiting for memory or the network card hardly moves.
Which of the two a firm's path is decides which server it should buy.

One Quant Book~13 (chapters~2 to~4 and~13) studied the processor, the memory hierarchy, the two-socket machine and the operating
system from the program's side. This chapter looks at the server as a thing bought and racked: how to choose its processor, what to
set in its firmware, what \omterm{def:nw:servers:oc}{overclocking} buys and costs, how it is managed, and how many of them fit in a cabinet.

\section{Choosing the processor}

\begin{definition}[Thermal design power]\label{def:nw:servers:tdp}
The \emph{thermal design power}\index{thermal design power} (TDP) of a processor is the sustained power, in watts, that its vendor
designs it to dissipate under load at its rated frequencies, and that the server's cooling and power supply must handle.
\end{definition}

Three numbers of a processor's specification matter for a hot path, and a fourth for the cabinet. The \emph{boost} clock is the
frequency one busy core can reach when the rest of the chip is quiet: AMD defines its EPYC processors' maximum boost as ``the maximum
frequency achievable by any single core''. The \emph{base} clock is what every core is guaranteed under full load. The cache size sets
how much of the order book and the strategy's state stays close to the core. And the \omterm{def:nw:servers:tdp}{thermal design power} sets how many servers a
cabinet's power allows.

\begin{proposition}[What a faster clock buys]\label{prop:nw:servers:clock}
Let a path take $t_{\mathrm{ref}}$ on a core clocked at $f_{\mathrm{ref}}$, of which a share $\phi$ is core-bound (it scales with the
clock) and the rest is bound by memory, the bus or the card. On a core at $f$ it takes
\[
  t(f) = t_{\mathrm{ref}}\Bigl(\phi\,\frac{f_{\mathrm{ref}}}{f} + 1 - \phi\Bigr),
\]
a speed-up that tends to $1/(1-\phi)$ however fast the clock.
\end{proposition}

\begin{proof}
The core-bound part is a number of cycles, which take $1/f$ each; the rest is a time fixed by the memory's and the devices' own
latencies. As $f \to \infty$ only $(1-\phi)\,t_{\mathrm{ref}}$ remains.
\end{proof}

\begin{example}[Book 13's path on three processors]\label{ex:nw:servers:book13}
One Quant Book~13's in-process stages (decode, ring, book, strategy, risk, gateway) had medians adding up to
\qty{1\,528}{\nano\second} on a laptop core whose maximum turbo is 4.8~GHz. Suppose 60\% of that time is core-bound. On a core at
5.0~GHz the path takes \qty{1\,491}{\nano\second}; at 4.1~GHz, \qty{1\,685}{\nano\second}; on a core held at a base clock of 2.7~GHz,
\qty{2\,241}{\nano\second}. The fastest clock of \cref{tab:nw:servers:candidates} saves 2.4\% over the reference; the slowest costs 47\%.
\end{example}

\begin{omfigure}
\begin{tikzpicture}
\begin{axis}[width=11.5cm,height=5.4cm,xmin=2.4,xmax=6.1,ymin=1100,ymax=3000,xlabel={core clock (GHz)},
  ylabel={in-process median (ns)},y tick label style={/pgf/number format/1000 sep={\,}},tick label style={font=\small},label style={font=\small},grid=major,grid style={gray!25},
  legend columns=4,legend style={font=\footnotesize,at={(0.5,-0.3)},anchor=north,draw=none,fill=none,
  /tikz/every even column/.append style={column sep=8pt}},table/col sep=comma]
\addplot[omThm,very thick] table[x=ghz,y=phi100]{figdata/networks/08-servers/model.csv};
\addplot[omMeth,very thick] table[x=ghz,y=phi60]{figdata/networks/08-servers/model.csv};
\addplot[omDef,very thick] table[x=ghz,y=phi30]{figdata/networks/08-servers/model.csv};
\addplot[only marks,mark=*,mark size=2pt,black] table[x=ghz,y=ns]{figdata/networks/08-servers/candidates.csv};
\legend{all core-bound,60\% core-bound,30\% core-bound,candidates (60\%)}
\end{axis}
\end{tikzpicture}

\omcaption{Model: Book~13's in-process path (\qty{1\,528}{\nano\second} at 4.8~GHz) moved to other clocks, for three shares of core-bound
time (\cref{prop:nw:servers:clock}); the dots are the candidates of \cref{tab:nw:servers:candidates} at their hot-core clock and 60\%.
The more the path waits for memory, the flatter the curve. Data: \texttt{fig\_servers.py}.}\label{fig:nw:servers:model}
\end{omfigure}

\begin{dated}{2026-09}{Server processors, from their vendors}\label{dat:nw:servers:parts}
AMD's EPYC 9175F: 16 cores, a base clock of 4.2~GHz, a maximum boost of up to 5~GHz, 512~MB of level-3 cache, a default TDP of
\qty{320}{\watt}, listed at 3\,703~USD in thousand-unit quantities (launched October 2024). AMD's EPYC 9755: 128 cores at 2.7~GHz base
and up to 4.1~GHz boost, 512~MB of level-3 cache, \qty{500}{\watt}, 10\,931~USD. Intel's Xeon Gold 6544Y: 16 cores at 3.6~GHz base and
4.1~GHz maximum turbo, 45~MB of cache, \qty{270}{\watt} (launched in the fourth quarter of 2023).
\end{dated}

\section{Firmware settings}

The firmware (BIOS) decides many things the operating system cannot change afterwards: how the processor saves power, whether each
core runs one thread or two, how memory is interleaved across sockets, and which events interrupt the processor behind the operating
system's back.

\begin{definition}[System management interrupt]\label{def:nw:servers:smi}
A \emph{system management interrupt}\index{system management interrupt} (SMI) is an interrupt handled by the platform's firmware, not
by the operating system, for functions such as temperature and fan control or legacy device emulation; it generally cannot be masked,
and while it runs the operating system and every program on that processor stop.
\end{definition}

An SMI is the one interrupt that core isolation (One Quant Book~13, chapter~13) cannot move away: Red Hat's real-time documentation warns
that a poorly written handler ``may consume many milliseconds'' that the operating system cannot preempt. Firms therefore ask their
server vendors which SMI sources can safely be turned off, count the ones that remain through the processor's counters, and treat an
unexplained pause of the whole machine as an SMI until shown otherwise.

\begin{method}[Firmware settings for a hot-path server]\label{met:nw:servers:bios}
\begin{enumerate}
\item Power profile for maximum performance; deep idle states (C-states) off on the hot cores, so a waiting core is never slow to wake.
\item A fixed turbo behaviour: either the maximum single-core boost with the other cores kept quiet, or a fixed all-core frequency, but
  not a clock that moves with load, which vendors describe as a source of jitter.
\item Simultaneous multithreading off, or its sibling threads left idle on the hot cores.
\item Memory at its maximum rated speed; interleaving across sockets off, so each socket's memory stays local.
\item Error-correction reporting and \omterm{def:nw:servers:smi}{system management interrupt} sources reduced to what the vendor states is safe.
\item The whole list recorded as data and audited, with the operating system's settings, before the server trades (the build).
\end{enumerate}
\end{method}

\section{Overclocked and liquid-cooled machines}

\begin{definition}[Overclocking]\label{def:nw:servers:oc}
\emph{Overclocking}\index{overclocking} is running a processor, or its memory, above its vendor's rated frequency, usually with a higher
supply voltage and stronger cooling, and outside the vendor's warranty.
\end{definition}

A faster clock at a higher voltage dissipates more power, roughly with the frequency times the square of the voltage, which is why the
specialist vendors pair \omterm{def:nw:servers:oc}{overclocking} with liquid cooling and oversized power supplies, and why they sell stability testing as part of the
product. What they sell is a processor that runs its cores at a frequency close to or above the vendor's single-core boost, all of
them, all day, without the clock changes that a load-dependent turbo brings.

\begin{dated}{2026-09}{Overclocked servers, from their vendors}\label{dat:nw:servers:oc}
Blackcore Technologies describes its ICON range as ``an overclocked, liquid-cooled server platform for Intel CPUs'', and its G3 SPR-M
server as a tuned Intel Xeon w7-2495X ``boasting 4.8GHz across all cores'' with 45~MB of cache and overclocked DDR5 memory, on a
custom 2\,000-watt redundant power supply. Business Systems International describes its ORION HF~310 as a single-socket overclocked server
whose design turns off unnecessary cores and locks the rest at up to 5.2~GHz, stating that ``turbo mode creates jitter, something that
traders want to avoid at all costs''.
\end{dated}

\begin{definition}[Burn-in test]\label{def:nw:servers:burnin}
A \emph{burn-in test}\index{burn-in test} runs a new or reconfigured server at full load, often for days, with stress workloads that
exercise the processor, memory and devices, and checks for errors, throttling and instability before the server is trusted with
trading.
\end{definition}

\section{Memory, reliability and management}

\begin{definition}[Error-correcting code memory]\label{def:nw:servers:ecc}
\emph{Error-correcting code memory}\index{error-correcting code memory} (ECC) stores extra check bits with each word, so that the memory
controller can correct a flipped bit and detect a double flip; server memory is almost always ECC, and overclocked memory makes its
reports more important, not less.
\end{definition}

\begin{definition}[Baseboard management controller]\label{def:nw:servers:bmc}
A \emph{baseboard management controller}\index{baseboard management controller} (BMC) is a small independent computer on the server's
motherboard, with its own network port, that lets an operator power the server on and off, read its sensors and logs, update its firmware
and use its console remotely, whatever the state of the main processor and operating system.
\end{definition}

A cage is often far from the firm's engineers. The BMC is how they reach a server that has hung, on the \omterm{def:nw:switches-and-layer-1-devices:oob}{out-of-band
management network} of chapter~2, and its sensors (temperatures, fan speeds, power) are the first sign that an overclocked machine is
drifting out of its tested envelope. It is also a computer with firmware of its own, on a network: it is patched and isolated like one.

\begin{omfigure}
\begin{tikzpicture}[font=\footnotesize,b/.style={draw,rounded corners=2pt,minimum height=0.6cm,align=center,inner sep=3pt},>=Stealth]
\node[b,fill=omDef!15,minimum width=2.6cm,minimum height=1.6cm] (s0) at (0,0) {socket 0\\ hot cores};
\node[b,fill=omDef!8,minimum width=2.6cm,minimum height=1.6cm] (s1) at (5.0,0) {socket 1\\ other work};
\foreach \k in {0,1,2,3} {\node[b,fill=gray!10,minimum width=0.5cm,font=\scriptsize] at (-1.0+0.65*\k,1.6) {M};}
\foreach \k in {0,1,2,3} {\node[b,fill=gray!10,minimum width=0.5cm,font=\scriptsize] at (4.0+0.65*\k,1.6) {M};}
\draw[<->,gray] (0,1.3) -- (s0.north);
\draw[<->,gray] (5.0,1.3) -- (s1.north);
\node[font=\scriptsize] at (-1.9,1.6) {memory};
\node[b,fill=omMeth!20] (nic) at (0,-2.0) {network card\\ (feed, orders)};
\node[b,fill=gray!8] (disk) at (5.0,-2.0) {storage, other cards};
\node[b,fill=omThm!12] (bmc) at (9.0,-0.9) {BMC};
\node[b,fill=gray!8] (mgmt) at (9.0,-2.4) {management\\ network};
\draw[<->,thick] (s0) -- node[above,font=\scriptsize]{interconnect} (s1);
\draw[<->,thick] (s0) -- node[right,font=\scriptsize]{PCIe} (nic);
\draw[<->,thick] (s1) -- (disk);
\draw[<->,gray] (bmc) -- (mgmt);
\draw[gray,dashed] (bmc.west) -- ++(-0.9,0) |- (s1.east);
\end{tikzpicture}

\omcaption{A two-socket trading server. The network card sits on socket~0's bus and the hot threads run on socket~0's cores with
socket~0's memory (One Quant Book~13, chapter~4); socket~1 takes everything else. The \omterm{def:nw:servers:bmc}{baseboard management controller} watches and
controls the machine from its own port on the management network.}\label{fig:nw:servers:anatomy}
\end{omfigure}

\section{Racks, power and heat}

\begin{definition}[Rack unit]\label{def:nw:servers:ru}
A \emph{rack unit}\index{rack unit} (U) is the standard height of equipment in a 19-inch rack, 1.75~inches (44.45~mm); a cabinet commonly
holds about 42 of them, and servers are 1U, 2U or taller.
\end{definition}

A colocation cabinet is sold with a power allowance in kilowatts (chapter~9), and power, not space, usually binds first. A server budgeted at
its processors' \omterm{def:nw:servers:tdp}{thermal design power} plus \qty{250}{\watt} for everything else draws \qty{570}{\watt} with one EPYC 9175F: seventeen fit in
\qty{10}{\kilo\watt}. The two-socket 128-core server draws \qty{1\,250}{\watt}: eight fit, with 2\,048 cores between them. The overclocked
server, budgeted at its \qty{2\,000}{\watt} supply because its vendor publishes no draw, fits five times.

\begin{table}[ht]
\centering\small
\begin{tabular}{lrrrrr}
\toprule
Candidate & Hot clock & Path at hot & Path at base & Servers per & Cores per \\
 & (GHz) & clock (ns) & clock (ns) & \qty{10}{\kilo\watt} & \qty{10}{\kilo\watt} \\
\midrule
AMD EPYC 9175F, one socket & 5.0 & 1\,491 & 1\,659 & 17 & 272 \\
Intel Xeon Gold 6544Y, one socket & 4.1 & 1\,685 & 1\,834 & 19 & 304 \\
AMD EPYC 9755, two sockets & 4.1 & 1\,685 & 2\,241 & 8 & 2\,048 \\
Overclocked w7-2495X (all cores) & 4.8 & 1\,528 & 1\,528 & 5 & 120 \\
\bottomrule
\end{tabular}
\caption{Candidates for a hot-path server, with Book~13's in-process path at 60\% core-bound (\cref{prop:nw:servers:clock}); the hot clock
is the single-core boost, or the vendor's all-core figure for the overclocked server. Servers per cabinet budget each at its processors'
TDP plus \qty{250}{\watt} (the overclocked one at its \qty{2\,000}{\watt} supply) and 42 \omterm{def:nw:servers:ru}{rack units}. Data: \texttt{nw\_servers.table()},
\cref{dat:nw:servers:parts,dat:nw:servers:oc}.}\label{tab:nw:servers:candidates}
\end{table}

\Cref{tab:nw:servers:candidates} makes the chapter's argument in one line each. The 16-core part with the highest boost is the fastest
for a path that runs on one or two cores, as long as the other cores stay quiet enough to let it boost. The overclocked machine gives up
a little of that peak for a clock that does not move, on every core, which matters when the hot path uses many cores at once. The
128-core part is the right server for research, simulation and the firm's other work (One Quant Book~15), and the wrong one for the hot
path.

\section{Tutorial: choosing and auditing a server}\label{tut:nw:servers:tut}

\begin{tutorial}
\textbf{Goal.} Predict the hot path's latency on candidate servers from Book~13's measurements, fit them into a cabinet, and audit a
server's firmware settings in the same report as its operating system's. \textbf{End state:} \cref{tab:nw:servers:candidates},
\cref{fig:nw:servers:model} and a combined audit report.
\begin{tutsteps}
\item \textbf{The reference.} \texttt{nw\_servers.ref\_ns()} reads Book~13's stage medians through \texttt{firm.wirepath} (1\,528~ns).
\item \textbf{The model.} \texttt{firm.serverspec.path\_ns} applies \cref{prop:nw:servers:clock}; \texttt{curve()} traces it and
  \texttt{python fig\_servers.py} writes the chart's data.
\item \textbf{The candidates.} \texttt{CANDIDATES} holds the ledger's rows; \texttt{table(phi, cabinet\_kw)} predicts each one's path and fits
  it into a cabinet.
\item \textbf{The audit.} \texttt{firmware\_audit(settings)} checks a server's firmware against \texttt{FIRMWARE} and returns
  \texttt{firm.tuneaudit} findings, so \texttt{firm\_tuneaudit.report} prints them with the host's own.
\end{tutsteps}
\textbf{What to change next.} Measure $\phi$ for Book~13's path by running it on the laptop's performance and efficiency cores (their maximum
clocks are 4.8 and 3.8~GHz) and fitting \cref{prop:nw:servers:clock}; recompute the table for a \qty{20}{\kilo\watt} cabinet.
\end{tutorial}

\omcode{code/firm/serverspec/firm_serverspec.py}{55}{75}{The latency model, a cabinet's fit, and the ranking of \texttt{firm.serverspec}.}\label{lst:nw:servers:model}

\section{Build: the server specification}\label{bld:nw:servers:serverspec}

\begin{build}
\bfield{Purpose} Server choice as data: candidates with their sources, a latency model tied to the firm's measured path, a cabinet fit, and
a firmware checklist audited with the host's tuning, for chapter~29's bill of materials.
\bfield{Interface} \texttt{firm\_serverspec}: \texttt{Candidate}, \texttt{CANDIDATES}, \texttt{path\_ns(t\_ref, phi, f, f\_ref)},
\texttt{speedup}, \texttt{per\_cabinet(c, kw, u)}, \texttt{rank}, \texttt{FIRMWARE}, \texttt{firmware\_audit(settings)} returning
\texttt{firm\_tuneaudit.Finding}s.
\bfield{Rules} Every candidate names its ledger row; a server's power budget is its processors' TDP plus a stated allowance, or a vendor's
figure; a setting that the server does not report is ``unknown'', never ``ok''.
\bfield{Acceptance tests} \texttt{code/firm/serverspec/tests/}: the model's limits (all core-bound, all memory-bound), a cabinet's fit by power
and by space, the ranking, and firmware findings reported through \texttt{firm.tuneaudit}.
\bfield{Stretch} Fit $\phi$ from runs at two clocks; add memory bandwidth to the model for a path whose memory-bound part depends on it.
\end{build}

\begin{omsources}
\item AMD, EPYC 9175F and EPYC 9755 product pages; Intel, Xeon Gold 6544Y and Core Ultra 7 155H specifications.
\item Blackcore Technologies, G3 SPR-M product page; Business Systems International, ORION HF~310 product page.
\item Red Hat, \emph{Red Hat Enterprise Linux for Real Time} reference guide (``System Management Interrupts'') and tuning guide (``Setting
BIOS parameters for system tuning'').
\end{omsources}

\section{Exercises}

\begin{exercise}[$\star$]\label{exo:nw:servers:1}
A path takes \qty{2}{\micro\second} at 3.0~GHz and is 80\% core-bound. How long does it take at 5.0~GHz, and what is its limit as the clock
grows?
\end{exercise}

\begin{exercise}[$\star$]\label{exo:nw:servers:2}
How many servers budgeted at \qty{570}{\watt} fit in a \qty{15}{\kilo\watt} cabinet of 42 units, if each takes one unit? And at
\qty{1\,250}{\watt} and two units?
\end{exercise}

\begin{exercise}[$\star$]\label{exo:nw:servers:3}
A server's boost clock is 5.0~GHz and its base clock 4.2~GHz. Which one should a latency budget use for the hot path, and when?
\end{exercise}

\begin{exercise}[$\star\star$]\label{exo:nw:servers:4}
Why can core isolation not protect a hot core from \omterm{def:nw:servers:smi}{system management interrupts}, and what can the firm do instead?
\end{exercise}

\begin{exercise}[$\star\star$]\label{exo:nw:servers:5}
In \cref{fig:nw:servers:model}, how much does raising the clock from 4.0 to 5.0~GHz save for each share of core-bound time?
\end{exercise}

\begin{exercise}[$\star\star$]\label{exo:nw:servers:6}
Why might a firm prefer an all-core overclocked machine to a higher single-core boost, even if its peak clock is lower?
\end{exercise}

\begin{exercise}[$\star\star\star$]\label{exo:nw:servers:7}
\emph{Coding.} With \texttt{firm.serverspec}, find the core-bound share below which the 16-core 5.0~GHz part saves less than 1\% over the
laptop reference, and redo \cref{tab:nw:servers:candidates} for a \qty{20}{\kilo\watt} cabinet.
\end{exercise}

\begin{exercise}[$\star\star\star$]\label{exo:nw:servers:8}
\emph{Find the flaw.} ``We moved the strategy from a 4.1~GHz server to a 5~GHz one and latency did not improve, so the benchmark was wrong.''
\end{exercise}

\section{Problem: Five Gigahertz or Fifty Cores}

\begin{problem}[{Weekend problem --- a clock, a cabinet and a budget}]\label{pb:nw:servers:1}
A firm's hot path takes \qty{1\,528}{\nano\second} at 4.8~GHz (One Quant Book~13's in-process stages). It is choosing between one-socket
servers with the EPYC 9175F (boost 5.0~GHz, base 4.2), two-socket servers with the EPYC 9755 (boost 4.1, base 2.7), and the overclocked
server of \cref{dat:nw:servers:oc} (4.8~GHz on all cores), for a \qty{10}{\kilo\watt} cabinet, with the chapter's power budgets.

\medskip\noindent\textbf{Part I --- The clock.}
\begin{enumerate}
\item At 60\% core-bound, what is the path on each server's hot clock?
\item And on each one's base clock?
\item What is the speed-up of the 9175F's boost over the 9755's base clock?
\item What is the largest speed-up any clock could give at 60\%?
\end{enumerate}

\medskip\noindent\textbf{Part II --- The cabinet.}
\begin{enumerate}[resume]
\item How many of each server fit in \qty{10}{\kilo\watt}?
\item How many cores does each give the cabinet?
\item Which binds, power or space, for each?
\item The firm needs 40 hot cores and 1\,000 other cores. How many cabinets of \qty{10}{\kilo\watt} does a mix of 9175F and 9755 servers need?
\end{enumerate}

\medskip\noindent\textbf{Part III --- The firmware.}
\begin{enumerate}[resume]
\item Which settings of the method would you check first on a new overclocked server, and why?
\item What does the \omterm{def:nw:servers:burnin}{burn-in test} protect against?
\item The BMC reports rising temperatures on one server. What do you do?
\item Why is SMT off on the hot cores?
\end{enumerate}

\medskip\noindent\textbf{Part IV --- The verdict.}
\begin{enumerate}[resume]
\item State the \emph{named result}: the hot path on the three servers at their hot clocks, and servers per cabinet.
\item Under what core-bound share does the choice of processor stop mattering by more than 5\%?
\item Where would you put the 128-core servers?
\item What would you measure before believing the model's $\phi$?
\item Why does the overclocked server's clock stability matter more than its peak?
\item What does the warranty question change in the purchase?
\item How would you present the choice to the firm's CTO in two sentences?
\item In one sentence: what decides whether a faster clock is worth buying?
\end{enumerate}
\end{problem}

\section{Interview questions}

\begin{interviewq}[$\star$ \iqroles{developer}]\label{iq:nw:servers:1}
Why do trading firms prefer processors with fewer, faster cores for the hot path?
\end{interviewq}

\begin{interviewq}[$\star\star$ \iqroles{developer}]\label{iq:nw:servers:2}
Which BIOS settings would you change on a new trading server, and why?
\end{interviewq}

\begin{interviewq}[$\star\star$ \iqroles{developer}]\label{iq:nw:servers:3}
What is a \omterm{def:nw:servers:smi}{system management interrupt} and why does it matter for latency?
\end{interviewq}

\begin{interviewq}[$\star\star$ \iqroles{developer, trader}]\label{iq:nw:servers:4}
Would you buy overclocked servers? What would you ask the vendor?
\end{interviewq}

\begin{interviewq}[$\star\star\star$ \iqroles{developer}]\label{iq:nw:servers:5}
How would you predict how much a faster processor will improve your path before buying it?
\end{interviewq}

\begin{interviewq}[$\star\star$ \iqroles{developer}]\label{iq:nw:servers:6}
Your colocation cabinet has \qty{10}{\kilo\watt}. How do you decide what goes in it?
\end{interviewq}
